\begin{claim}{D.59}{Proved}{The negative cusp section and its fixed homology line}
Let $b$ be a negative involution in the full
once-punctured-torus loop-outer group, let
$F=\widehat F(x,y)$ and $c=[x,y]$, and choose a
cusp reflection $v$ in its centre-free extension $E$.
Thus
\[
 E=F\rtimes\langle v\rangle,\qquad
 v^2=1,\qquad vcv^{-1}=c^{-1}.
\]
Put $\alpha=\operatorname{Inn}(v)|_F$,
$C=C_F(v)$, and
\[
 M=H_1(F,\widehat\Z),\qquad
 M^+=\ker(\alpha_*-I:M\to M).
\]
Then the cusp section is specifically the class with
\[
             H^1(C,\F_2)=0,\qquad C\subseteq\overline{[F,F]}.
                                                        \tag{D.59a}
\]
For every actual transporter $g$ supplied by D.58,
\[
                    gvg^{-1}=cv,
                                                        \tag{D.59b}
\]
its homology class is a generator of the full fixed
line:
\[
 \widehat\Z\xrightarrow{\ \cong\ }M^+,\qquad
                         t\longmapsto t[g].
                                                        \tag{D.59c}
\]
This line is a direct summand of $M$.

More precisely, let $m\geq3$ be odd and let
$\chi:F\to C_m$ be any surjective continuous
$\alpha$-invariant character. Extend it to
$\widetilde\chi:E\to C_m$ by
$\widetilde\chi(v)=0$, and put
$F_m=\ker\chi$, $E_m=\ker\widetilde\chi$.
Then
\[
 C\subset F_m,\qquad
 \chi(g)\text{ generates }C_m,\qquad
 E_m=F_m\rtimes\langle v\rangle .
                                                        \tag{D.59d}
\]
The $m$ cusp-dihedral edges form one cycle on the
$m$ $F_m$-section classes above the original cusp
class. Their fixed groups are conjugates of the
same full group $C$.
The complete section counts in this cover are
\[
\begin{array}{c|c|c}
 \alpha_*\bmod2&
 \text{cusp classes in }E_m&
 \text{other classes in }E_m\\ \hline
 I&m\text{, each with }H^1(-,\F_2)=0&
 1\text{, with }\dim H^1(-,\F_2)=1\\
 \ne I&m\text{, each with }H^1(-,\F_2)=0&0 .
\end{array}                                             \tag{D.59e}
\]
In the first row, if $D=C_F(w)$ is the fixed group
of the other full section class, its actual
homology-image map is surjective:
\[
                         D\longrightarrow M^+ .
                                                        \tag{D.59f}
\]
These statements specify full subgroup cohomology
and the indicated ambient homology maps. They make
no identification of $C$ with the trivial group,
of $D$ with a procyclic group, or of $g$ with a member
of a full free profinite basis.
\end{claim}

\begin{proof}
D.43 gives the global multiplier $-1$.
D.58 gives the stated actual reflection, the
transporter (D.59b), and exactly one or two full
section classes according to the two rows above.
Their normalizers are the actual direct products
$C_F(S)\times S$.
Write $A_2$ for the action on $H_1(F,\Z_2)$.

\emph{The odd-degree cover and its boundary injection.}
The displayed extension of $\chi$ is a homomorphism
because $\chi\alpha=\chi$.
Its kernel is precisely the split extension in
(D.59d). Every order-two element of $E$ lies in
that kernel, since a homomorphism to the odd cyclic
group kills such an element.

The restriction of $\chi$ to the dense discrete
free group is onto. Its degree-$m$ covering of
the punctured torus has $m$ degree-one boundary
lifts, because $\chi(c)=0$.
Filling those boundaries gives the corresponding
unramified degree-$m$ cover of the compact torus.
Consequently its type is $(1,m)$, its completed
group is exactly $F_m$, and
$F_m\cong\widehat F_{m+1}$.
Choose a discrete word $f$ with $\chi(f)=1$, and put
\[
 c_i=f^icf^{-i},\qquad
 v_i=f^ivf^{-i},\qquad 0\leq i<m .
                                                        \tag{D.59g}
\]
Each $v_i$ belongs to $E_m$, has order two, and
inverts $c_i$.

Here is the exact homology comparison required
for the section argument.
Let $\Delta=F/F_m=C_m$ be the deck group and let
$P=H_1(F_m,\Z_2)$. The boundary sequence is
\[
 0\longrightarrow B\longrightarrow P
       \longrightarrow H_1(\overline S_m,\Z_2)
       \longrightarrow0,\qquad
 B=\Z_2^m/\Z_2(1,\ldots,1).
                                                        \tag{D.59h}
\]
The image of the $i$th standard generator of $B$
is $[c_i]$. Thus these classes have their sole
relation $\sum_i[c_i]=0$.
Deck transformations act trivially on the compact
torus homology. Since $m$ is odd, the actual
idempotent
\[
                    e=\frac1m\sum_{\delta\in\Delta}\delta
                                                        \tag{D.59i}
\]
is defined on $P$. It annihilates $B$ and induces
the identity on the last term of (D.59h).
It follows that
\[
             P=eP\oplus(1-e)P,\qquad (1-e)P=B .
                                                        \tag{D.59j}
\]
The covering map identifies $eP$ with the base
torus homology: on integral homology its image
is the index-$m$ covering lattice, and that
inclusion is an isomorphism over $\Z_2$.

The automorphism $\alpha$ induces the identity
on $\Delta$. Its action on $P$ therefore commutes
with the deck action; a discrepancy between two
chosen lifts is an inner automorphism of $F_m$
and has trivial action on $P$.
It preserves every individual boundary class
with multiplier $-1$. Explicitly,
$\alpha(f^i)f^{-i}\in F_m$ conjugates $c_i^{-1}$
to $\alpha(c_i)$. Thus $\alpha=-I$ on $B$.
Its action on $eP$ is the actual base matrix $A_2$,
through the covering homology map just specified.

Reduction modulo two makes $\alpha$ the identity
on $B/2B$. The abelianization of the split extension
gives the exact comparison
\[
 H_1(E_m,\F_2)
   \cong \bigl(H_1(F_m,\F_2)\bigr)_{\alpha}
                    \oplus\F_2[v].
                                                        \tag{D.59k}
\]
The direct summand $B/2B$ survives injectively
in these coinvariants by (D.59j).
In particular the images
$d_i=[c_i]\in H_1(E_m,\F_2)$ still have precisely
their one total-sum relation. Their span has
dimension $m-1$.

\emph{The actual section graph and its transporter labels.}
Let $h=|\chi(C)|$.
The $F_m$-conjugacy classes lying above the
original full cusp-section class are exactly
the vertices
\[
                              C_m/\chi(C).
                                                        \tag{D.59l}
\]
Indeed $f^ivf^{-i}$ and $f^jvf^{-j}$ are
$F_m$-conjugate precisely when $i-j\in\chi(C)$:
a conjugator $n\in F_m$ gives
$f^{-j}nf^i\in C$, and the converse follows by
reversing that formula.
All original order-two sections lie in $E_m$,
so this is the full class splitting, not a
selection of section images in a finite quotient.

For each $i$, take the edge with endpoints
$\langle v_i\rangle$ and $\langle c_i v_i\rangle$.
Both endpoints are actual order-two sections.
Equation (D.59b) identifies the second endpoint
with the conjugate by $f^ig$ of $\langle v\rangle$.
Thus, with $\delta=\chi(g)$, this edge has the
exact vertex labels
\[
              i+\chi(C)\quad\text{and}\quad
                         i+\delta+\chi(C).
                                                        \tag{D.59m}
\]
There are $m$ labelled edges and $m/h$ vertices.
Assign to a vertex the homology class in
$H_1(E_m,\F_2)$ of its nonidentity section element.
This is independent of the conjugate representative.
The sum of the two endpoint values of the $i$th
edge is exactly $d_i$, by the group equality
$c_i v_i=c_i\cdot v_i$.

For a finite graph with $r$ vertices and $k$
connected components, its edge differences span
a subspace of dimension at most $r-k$ in any
vector space of vertex values.
One proves this by choosing one vertex in each
component and expressing every edge difference
in the differences from those chosen vertices.
Apply this to the actual values above. By (D.59k),
\[
                         m-1\leq m/h-k .
                                                        \tag{D.59n}
\]
Since $h\geq1$ and $k\geq1$, equality is forced:
$h=1$ and $k=1$.
The edges now join $i$ to $i+\delta$ in $C_m$.
Connectedness says exactly that $\delta$ generates
$C_m$. Since $m$ is odd and at least three,
these edges form a single $m$-cycle.
This proves all assertions in (D.59d) and the
cusp part of the graph assertion.
Because $C\subset F_m$, the fixed group of each
such section is the corresponding conjugate of
the same full group $C$.

\emph{Which full section is the cusp section.}
We compute the high-degree cohomology of $E_m$
using the actual module (D.59j).
Over $\F_2$ its dual consists of $m-1$ trivial
boundary summands and the dual base module.
When $A_2\bmod2=I$ this last module is two
trivial summands; otherwise it is the regular
two-dimensional $\F_2[C_2]$-module.
The two-row Hochschild--Serre sequence has no
nonzero differential into the quotient row:
the actual section makes inflation from $C_2$
injective in every degree. There are no other
possible differentials.
The periodic cyclic resolution therefore gives,
for every $n\geq2$,
\[
 \dim_{\F_2}H^n(E_m,\F_2)=
 \begin{cases}
   m+2,&A_2\bmod2=I,\\
   m,&A_2\bmod2\ne I.
 \end{cases}                                           \tag{D.59o}
\]

The full-profinite Brown restriction theorem
applies with the open subgroup $F_m$ of
$2$-cohomological dimension one
\cite[Theorem~2 and \S1.1]{A-KochloukovaZalesskii2007}.
Every nontrivial finite subgroup has order two.
For a section $T$, its normalizer in $E_m$ is
$C_{F_m}(T)\times T$, with
$\operatorname{cd}_2C_{F_m}(T)\leq1$.
The finite source in (D.59o) makes the dense
Brown restriction an isomorphism onto its
finite product, so each section class contributes
\[
                     1+\dim H^1(C_{F_m}(T),\F_2)
                                                        \tag{D.59p}
\]
to (D.59o).
Conjugacy by $E_m$ and by $F_m$ agrees: multiply
an $E_m$-conjugator by the element of its source
section with the same quotient label.

D.58 already gives $H^1(C,\F_2)=0$ when
$A_2\bmod2\ne I$.
In the other case its dimension is zero or one.
Choose an invariant surjective character to
$C_3$: the integral involution on $\Z_3^2$ has
distinct eigenvalues $1,-1$, so projection onto
its fixed direct summand supplies this character.
The graph just proved gives three cusp classes
with the same group $C$. A dimension-one $H^1(C,\F_2)$
would make them contribute six to (D.59o),
whose value is five. Thus the dimension is zero
also in this case, proving the first assertion
of (D.59a).

If $A_2\bmod2\ne I$, D.58 has no other full
section class, so the $m$ cusp classes exhaust
$E_m$. If $A_2\bmod2=I$, let $D=C_F(w)$ be the
fixed group of the other full class.
As observed above, $w\in E_m$.
The same exact class-splitting formula gives
$m/|\chi(D)|$ classes, each with fixed group
conjugate to $D\cap F_m$. Subtracting the $m$
cusp contributions from (D.59o) yields
\[
 \frac{m}{|\chi(D)|}
       \bigl(1+\dim H^1(D\cap F_m,\F_2)\bigr)=2 .
                                                        \tag{D.59q}
\]
The first factor is odd. It is therefore one,
and the cohomology dimension is one.
This proves (D.59e), as well as
$\chi(D)=C_m$ for every such odd character.

\emph{The full ambient homology images.}
Every element of $C$ has homology in $M^+$.
At an odd prime $p$, the two idempotents
$(I\pm A_p)/2$ split $\Z_p^2$ into its rank-one
fixed and negative summands.
For every $r\geq1$, projection onto a chosen
generator of the fixed summand followed by
reduction modulo $p^r$ is an actual invariant
surjective character $F\to C_{p^r}$.
Equation (D.59d) makes the image of $C$ vanish
under every one of these maps, and hence in the
entire fixed $\Z_p$-line.
At the prime two, a nonzero compact image of
$C\to M^+_2\cong\Z_2$ would be $2^r\Z_2$ for
some $r$, and its quotient by twice itself would
give a nonzero continuous character $C\to\F_2$.
The already proved $H^1(C,\F_2)=0$ excludes that.
Thus $C$ has zero image in every primary factor
of $M$, which is exactly
$C\subseteq\ker(F\to M)=\overline{[F,F]}$.

From $c=g\alpha(g)^{-1}$, the homology of $g$
belongs to $M^+$. At every odd prime the same
characters and the generator assertion of
(D.59d) say that $[g]$ generates the fixed
$\Z_p$-line. We check the dyadic factor directly.

In the maximal pro-two group, let $D_2,D_3$
be the second and third Zassenhaus terms.
The quotient $D_2/D_3$ is the three-dimensional
$\F_2$-space with basis
\[
                        X^2,\quad Y^2,\quad Z=[X,Y].
\]
For completeness this is the central subgroup of
the finite characteristic quotient with
coordinates $(a,b,k)\in(\Z/4)^2\times\F_2$ and
product
\[
 (a,b,k)(a',b',k')
     =(a+a',\,b+b',\,k+k'+ba').
                                                        \tag{D.59r}
\]
The generators are $(1,0,0),(0,1,0)$ and
their commutator is $(0,0,1)$.
The subgroup with $a,b$ even has the stated basis.
Equivalently its kernel is the closed subgroup
generated by fourth powers, squares of commutators,
and triple commutators, so the quotient is
characteristic.

The induced action on this degree-two space depends
only on $\overline A=A_2\bmod2$. A vector
$aX+bY$ has square
$aX^2+bY^2+abZ$, and $\overline A$ fixes $Z$.
If $\overline A=I$, the degree-two action is
the identity. Otherwise an actual linear basis
change over $\F_2$ gives
$X\mapsto X$, $Y\mapsto Y+X$.
It preserves $Z$ since every invertible
two-dimensional matrix over $\F_2$ has determinant
one. In these coordinates
\[
 \operatorname{im}(I+\overline A|_{D_2/D_3})
                        =\langle X^2+Z\rangle .
                                                        \tag{D.59s}
\]
In either case $Z$ is not in that image.
If $[g]$ vanished modulo two, its pro-two image
would lie in $D_2$, and the exact equation
$c=g\alpha(g)^{-1}$ would give
$Z=(I+\overline A)[g]_2$ in $D_2/D_3$.
This contradicts (D.59s).
Thus $[g]\bmod2\ne0$.

The kernel of $A_2-I$ on $\Z_2^2$ is a saturated
rank-one submodule: if $2z$ lies in it, then
$2(A_2-I)z=0$ in a torsion-free module.
It is therefore generated by a primitive vector,
and the nonzero reduction of $[g]$ makes its
coefficient in that generator a unit.
Together with the odd-primary conclusions, this
proves (D.59c).
For each prime choose a complementary basis vector
to $[g]_p$. Taking their product gives an actual
decomposition $M=\widehat\Z[g]\oplus\widehat\Z e$
and a continuous homology retraction
$M\to\widehat\Z$ sending $[g]$ to one.

Finally suppose $A_2\bmod2=I$.
Equation (D.59q) makes the ambient image of $D$
surject onto every odd-primary fixed line.
At two, use the actual positive character
$\varphi_+$ and its reduction $a$ in D.58e--D.58j.
The nonzero nilpotent class $a/u$ restricts under
the full Brown map to the one-dimensional
$H^1(D,\F_2)$ factor; the cusp factor has no such
summand by (D.59a). Hence $a|_D\ne0$.
The integral idempotent decomposition of D.58
makes $\varphi_+$ primitive on the fixed dyadic
line. Since the ambient image of $D$ lies in
that line, its image contains a dyadic unit
and is all of the line.
The total image is a compact subgroup of
$M^+\cong\widehat\Z$ surjecting onto every primary
factor. Its image in each finite cyclic quotient
is therefore the full quotient; compactness
proves surjectivity onto $\widehat\Z$.
This establishes (D.59f).

The executable
\path{code/part_c/torus_cusp_cycle_d59.py}
checks the exact finite incidence ranks, boundary
relations, cyclic-module dimensions and the
degree-two quotient calculation. The full group
and topology comparisons, including the Brown
restriction and all ambient homology maps, are
proved above.
\end{proof}
